\documentclass[11pt]{article}

\usepackage[margin=1in]{geometry}
\usepackage[T1]{fontenc}
\usepackage{lmodern}
\usepackage{fontspec}
\setmainfont{Calibri}
\usepackage{amsmath,amssymb}
\usepackage{graphicx}
\usepackage{enumitem}
\usepackage{url}
\usepackage[strings]{underscore}
\usepackage[hidelinks]{hyperref}

\newcommand{\gmfig}[2]{%
  \IfFileExists{gruggle_testsuite/#1.pdf}{\includegraphics[width=#2]{gruggle_testsuite/#1.pdf}}{%
    \texttt{[Missing PDF figure gruggle\_testsuite/#1.pdf]}}}

\title{Gruggle: Lightweight Graphviz DOT Manipulation}
\author{Sambuddha Majumder, Partha P Chakrabarti}
\date{\today}

\begin{document}
\maketitle

\begin{abstract}
Gruggle is a small Node.js tool for ingesting, merging, inspecting, and
lightly manipulating Graphviz DOT graphs from the command line.
It combines a permissive DOT parser (supporting comments, ports,
attributes, subgraphs) with a composable subcommand pipeline that can
add nodes, add edges (optionally random), prune isolated nodes, perform
path-based keep/remove filtering, and export to SVG via Graphviz \texttt{dot}.
\end{abstract}

\section{Command Line and Entry Point}

The main executable is \texttt{gruggle.js}.  It accepts zero or more
input files using repeated \texttt{-i} flags.  All parsed graphs are
merged; if no inputs are given an empty digraph is used.
Style switches:
\begin{itemize}[leftmargin=*]
  \item \texttt{--curved} (default): curved spline edges
  \item \texttt{--ortho} or \texttt{--rectilinear}: rectilinear edges (\texttt{splines=ortho})
\end{itemize}
After parsing/merging, subcommands are applied in order.  Examples:
\begin{verbatim}
node gruggle.js --ortho -i g1.dot -i g2.dot set-name merged to-svg merged.svg
node gruggle.js chain "add-node nx:1:100" "add-edge 50 nx.* nx.*" \
                     "keep-paths nx1 nx5" "to-svg pruned.svg"
\end{verbatim}

\section{DOT Parsing and Merging}

The parser handles both \texttt{graph}/\texttt{digraph}, optional
\texttt{strict}, named/unnamed graphs, subgraphs, edge ops
\texttt{->}/\texttt{--}, port/compass points, bracketed attribute
lists, quoted IDs, HTML-like IDs, numbers, comments (\texttt{//},
\texttt{/* */}, \texttt{\#}), and default attribute statements
(\texttt{graph}, \texttt{node}, \texttt{edge}).  When merging, later
files' attributes override earlier ones on conflicts; edges are
deduplicated by endpoints/ports/direction with attribute merge.

Default emitted style: box-shaped, rounded, light-blue nodes and
muted-blue edges; orthogonal vs.\ curved splines are controlled by the
CLI switch above.

\section{Subcommands}

Subcommands operate on the in-memory graph and can be composed with
the \texttt{chain} meta-command.  Tokens inside a chain are parsed with
a simple shell-like splitter (quotes supported).
\begin{itemize}[leftmargin=*]
  \item \textbf{set-name} \texttt{<name>}: rename the graph.
  \item \textbf{add-nodes} \texttt{<id>|<prefix:start:end>}: add a
        single node or a numeric range (e.g.\ \texttt{nx:1:3}
        $\rightarrow$ \texttt{nx1}, \texttt{nx2}, \texttt{nx3}).
  \item \textbf{add-edges} \texttt{<label> <fromRe> <toRe>}: add edges
        for all pairs whose IDs match the regexes, using the current
        graph direction.  Attribute \texttt{label} is set when
        provided and non-empty.
  \item \textbf{add-edges} \texttt{<count> <fromRe> <toRe> [label]}:
        pick up to \texttt{count} random distinct pairs from the regex
        matches and add those edges (optional label).
  \item \textbf{add-chain} \texttt{<root> <n>}: add nodes
        \texttt{root\_1}.. \texttt{root\_n} linked in a chain.
  \item \textbf{add-grid} \texttt{<root> <rows> <cols>}: add a
        \texttt{rows} x \texttt{cols} grid with nodes
        \texttt{root\_i\_j} and directional grid edges (top-to-bottom,
        left-to-right).
  \item \textbf{delete-isolated-nodes}: remove nodes with no incident
        edges.
  \item \textbf{remove-nodes} \texttt{<regex>}: delete nodes matching
        the regex and their incident edges.
  \item \textbf{remove-edges-between-nodes} \texttt{<fromRe> <toRe>}:
        delete edges whose tail/head IDs match the regex pair.
  \item \textbf{add-labels-to-edges-on-paths} \texttt{<fromRe> <toRe> <label>}:
        add the given label (appended to any existing comma-separated
        labels) on every edge lying on a path from sources matching
        \texttt{fromRe} to targets matching \texttt{toRe}.
  \item \textbf{remove-labels-from-edges-on-paths} \texttt{<fromRe> <toRe> <label>}:
        remove the given label from every edge lying on such paths; if
        no labels remain, the attribute is removed.
  \item \textbf{select-nodes} / \textbf{deselect-nodes} \texttt{<regex>}:
        mark/unmark nodes whose IDs match the regex.
  \item \textbf{select-edges} / \textbf{deselect-edges}:
        either \texttt{<labelRegex>} (matches any label component) or
        \texttt{<fromRe> <toRe>} (matches endpoints) to mark/unmark edges.
  \item \textbf{select-nodes-label} \texttt{<labelRegex>}: select nodes
        whose label attribute has a matching component (comma-separated).
  \item \textbf{delete-selected}: delete selected nodes (and their
        incident edges) and selected edges.
  \item \textbf{set-node-attr-on-selected} \texttt{<attr> <value>} /
        \textbf{set-edge-attr-on-selected} \texttt{<attr> <value>}:
        set a DOT attribute on selected nodes/edges.
  \item \textbf{unset-node-attr-on-selected} \texttt{<attr>} /
        \textbf{unset-edge-attr-on-selected} \texttt{<attr>}:
        remove a DOT attribute on selected nodes/edges.
  \item \textbf{clear-visual-attrs-on-selected}: clear common visual
        attributes (color, fillcolor, fontcolor, style, shape,
        penwidth, arrowsize, fontsize) on selected nodes/edges.
  \item \textbf{add-labels-to-selected} \texttt{<label>}: append the
        label (comma-separated) to all selected nodes and edges.
  \item \textbf{remove-labels-from-selected} \texttt{<label>}: remove
        the label from selected nodes/edges; drops the attribute if empty.
  \item \textbf{keep-paths} \texttt{<fromRe> <toRe>}: mark all nodes
        and edges that lie on at least one path from any source
        matching \texttt{fromRe} to any target matching \texttt{toRe}.
        Marks set \texttt{keep=true} and clear \texttt{remove}.
  \item \textbf{remove-paths} \texttt{<fromRe> <toRe>}: similarly mark
        nodes/edges on such paths, setting \texttt{remove=true} and
        clearing \texttt{keep}.
  \item \textbf{keep-path-edges} \texttt{<fromRe> <toRe>}: mark only
        the edges (not nodes) that lie on such paths with
        \texttt{keep=true}.
  \item \textbf{remove-path-edges} \texttt{<fromRe> <toRe>}: mark only
        the edges (not nodes) on such paths with
        \texttt{remove=true}.
  \item \textbf{filter-to-paths} \texttt{<fromRe> <toRe>}: keep only
        nodes/edges on such paths (mark all others as remove).
  \item \textbf{select-nodes-on-paths} \texttt{<fromRe> <toRe>}:
        select all nodes on such paths (for later bulk ops).
  \item \textbf{select-edges-on-paths} \texttt{<fromRe> <toRe>}:
        select all edges on such paths.
  \item \textbf{find-path} \texttt{<regex> <label>}: run greggle
        \texttt{find-path} to annotate one matching path's edges with
        the given label using greggle-compatible regular-path
        semantics. The regex may include lifted node predicates such
        as \texttt{P@} and \texttt{@P}.
  \item \textbf{select-neighbours} \texttt{[in|out|both] <hops>}:
        expand selection by BFS over neighbours (default both,1).
  \item \textbf{select-nodes-degree} \texttt{<lt|le|eq|ge|gt> <k>}:
        select nodes by total degree (in+out for digraphs).
  \item \textbf{select-nodes-indegree} \texttt{<lt|le|eq|ge|gt> <k>}:
        select nodes by indegree (for undirected graphs, each edge
        counts as both in and out).
  \item \textbf{select-nodes-outdegree} \texttt{<lt|le|eq|ge|gt> <k>}:
        select nodes by outdegree (for undirected graphs, each edge
        counts as both in and out).
  \item \textbf{select-component-of-nodes} \texttt{<regex>}: select
        nodes/edges in components of matching nodes.
  \item \textbf{select-components-with-selected}: select nodes/edges in
        any component that already contains a selected node.
  \item \textbf{clear-selection}: clear selection on all nodes/edges.
  \item \textbf{invert-selection-nodes}: flip selection flags on nodes.
  \item \textbf{invert-selection-edges}: flip selection flags on edges.
  \item \textbf{select-neighbours} \texttt{[in|out|both] <hops>}:
        expand selection by BFS over neighbours (default both,1).
  \item \textbf{select-nodes-degree} \texttt{<lt|le|eq|ge|gt> <k>}:
        select nodes by total degree (in+out for digraphs).
  \item \textbf{chain} \texttt{"cmd1 ..."} \texttt{"cmd2 ..."}: run
        each nested command in order. Inside a chain, \texttt{-f
        <file>} splices commands line-by-line from the file
        (comments/blank lines skipped) at that point in sequence.
  \item \textbf{to-svg} \texttt{<filename>}: pipe the current DOT
        through \texttt{dot -Tsvg -o <filename>} (assumes \texttt{dot}
        is on \texttt{PATH}).
\end{itemize}

\paragraph{Keep/remove semantics.}
Elements carry in-memory \emph{keep} and \emph{remove} flags that flip
with the operations above (keep clears remove, remove clears keep).
When emitting the final graph:
\begin{itemize}[leftmargin=*]
  \item If any element is marked keep, only kept elements are emitted.
  \item Otherwise, elements marked remove are suppressed; the rest are
        emitted.
\end{itemize}
No flags are written into DOT attributes.

\section{Output}

Default output is DOT on stdout, reflecting merged defaults and the
current subcommand state.  Using \texttt{to-svg} writes an SVG file
via Graphviz with the same styling preamble (boxy light-blue nodes,
blue edges, curved or orthogonal per CLI switch).

\section{Example Workflow}
\begin{verbatim}
# Merge two graphs, add 20 random edges among nx* nodes, retain only
# paths from nx1 to nx5, emit both DOT and SVG with orthogonal edges.
node gruggle.js --ortho -i a.dot -i b.dot \
  chain "add-edge 20 nx.* nx.*" "keep-paths nx1 nx5" "to-svg pruned.svg" \
  > pruned.dot
\end{verbatim}

\section{Illustrative Examples}

\subsection{Ring/Chain/Grid Builders}
Figure~\ref{fig:chain-ring} and Figure~\ref{fig:grid} illustrate the synthetic builders discussed here.
\begin{figure}[h]
  \centering
  \gmfig{add_chain.golden}{0.4\textwidth}
  \gmfig{add_ring.golden}{0.4\textwidth}
  \caption{Left: chain of four nodes. Right: ring of three nodes.}
  \label{fig:chain-ring}
\end{figure}
\begin{figure}[h]
  \centering
  \gmfig{add_grid.golden}{0.5\textwidth}
  \caption{2x2 grid built with \texttt{add-grid g 2 2}.}
  \label{fig:grid}
\end{figure}
\begin{verbatim}
# Chain of four nodes c_1 -> c_2 -> c_3 -> c_4
node gruggle.js chain "add-chain c 4" > chain.dot

# Ring of three nodes r_0 -> r_1 -> r_2 -> r_0
node gruggle.js chain "add-ring r 3 0" > ring.dot

# 2x2 grid g_1_1, g_1_2, g_2_1, g_2_2 with edges down/right
node gruggle.js chain "add-grid g 2 2" > grid.dot
\end{verbatim}

\subsection{Random Selection and Labeling}
Figure~\ref{fig:rnd} shows random node/edge selection annotated with \texttt{rnd}.
\begin{figure}[h]
  \centering
  \gmfig{select_random_nodes_edges.golden}{0.45\textwidth}
  \caption{Randomly selected nodes/edges labeled \texttt{rnd}.}
  \label{fig:rnd}
\end{figure}
\begin{verbatim}
# Select up to 10 nodes and edges at random, mark them 'rnd'
node gruggle.js -i sample1.dot chain \
  "select-random-nodes 10 .*" \
  "select-random-edges 10 .* .*" \
  "add-labels-to-selected rnd" > random.dot
\end{verbatim}
Resulting edges/nodes carry labels like \texttt{name:rnd} or \texttt{label="x,rnd"}.

\subsection{Path Highlighting via Greggle (find-path)}
Figure~\ref{fig:pathlab} shows greggle-driven path annotation.
\begin{figure}[h]
  \centering
  \gmfig{find_path.golden}{0.5\textwidth}
  \caption{Edges on one matching path annotated with \texttt{pathlab}.}
   \label{fig:pathlab}
\end{figure}
\begin{verbatim}
# Highlight one path whose edge labels match (concat x y) with 'pathlab'
node gruggle.js -i sample1.dot chain \
  "find-path (concat x y) pathlab" > path_highlighted.dot
\end{verbatim}
Edges on the witness path gain \texttt{pathlab} in their label (e.g., \texttt{label="x,pathlab"}).
The same command also supports lifted node predicates inside the regex.
For example, \texttt{(concat a (all-of a @TARGET))} selects a
two-edge witness whose second edge is labelled \texttt{a} and whose
sink node label matches \texttt{TARGET}.

\subsection{Component-Based Selection}
Figure~\ref{fig:components} visualizes component selection from regex and from existing selection.
\begin{figure}[h]
  \centering
  \gmfig{select_component_of_nodes.golden}{0.45\textwidth}
  \gmfig{select_components_with_selected.golden}{0.45\textwidth}
  \caption{Selecting components by node regex (left) and from existing selection (right).}
  \label{fig:components}
\end{figure}
\begin{verbatim}
# Select components containing nodes matching x*
node gruggle.js -i comps.dot chain \
  "select-component-of-nodes x.*" \
  "add-labels-to-selected compA" > comps_marked.dot

# Select components containing already-selected nodes
node gruggle.js -i comps.dot chain \
  "select-nodes y2" \
  "select-components-with-selected" \
  "add-labels-to-selected compB" > comps2.dot
\end{verbatim}
Nodes and internal edges in those components are labeled \texttt{name:compA} or \texttt{name:compB}.

\subsection{Selection Utilities}
Figure~\ref{fig:visattrs} shows styling and clearing of selected elements.
\begin{figure}[h]
  \centering
  \gmfig{visual_attrs_on_selected.golden}{0.45\textwidth}
  \caption{Visual attrs set then cleared on selected nodes/edges.}
  \label{fig:visattrs}
\end{figure}
\begin{verbatim}
# Label-selected nodes/edges, then clear visual styling
node gruggle.js -i sample1.dot chain \
  "select-nodes s1" \
  "select-edges x" \
  "set-node-attr-on-selected fillcolor yellow" \
  "set-edge-attr-on-selected color red" \
  "clear-visual-attrs-on-selected" > styled.dot
\end{verbatim}

\subsection{Degree and Neighborhood Selection}
Figure~\ref{fig:deg-neigh} illustrates degree-based selection and neighbour expansion before the corresponding command snippets.
\begin{figure}[h]
  \centering
  \gmfig{select_nodes_degree.golden}{0.45\textwidth}
  \gmfig{select_neighbours.golden}{0.45\textwidth}
  \caption{Left: nodes with degree $\ge 3$ labeled \texttt{degsel}. Right: 2-hop neighbour expansion labeled \texttt{neigh}.}
  \label{fig:deg-neigh}
\end{figure}
\begin{verbatim}
# Select nodes with degree >= 3
node gruggle.js -i sample1.dot chain \
  "select-nodes-degree ge 3" \
  "add-labels-to-selected degsel" > deg.dot

# BFS expand selection by 2 hops (undirected sense)
node gruggle.js -i sample1.dot chain \
  "select-nodes s1" \
  "select-neighbours both 2" \
  "add-labels-to-selected neigh" > neigh.dot
\end{verbatim}

\subsection{Directed Indegree/Outdegree Selection}
Figure~\ref{fig:indeg-outdeg} shows indegree and outdegree filtering on a directed sample graph.
\begin{figure}[h]
  \centering
  \gmfig{select_nodes_indegree.golden}{0.45\textwidth}
  \gmfig{select_nodes_outdegree.golden}{0.45\textwidth}
  \caption{Left: nodes with indegree $\ge 2$ labeled \texttt{indegsel}. Right: nodes with outdegree $=0$ labeled \texttt{outdegsel}. Undirected edges count toward both in/out.}
  \label{fig:indeg-outdeg}
\end{figure}
\begin{verbatim}
# Indegree >= 2 (counts incoming arcs; undirected edges count as both)
node gruggle.js -i sample1.dot chain \
  "select-nodes-indegree ge 2" \
  "add-labels-to-selected indegsel" > indeg.dot

# Outdegree == 0 (no outgoing arcs; undirected edges count as both)
node gruggle.js -i sample1.dot chain \
  "select-nodes-outdegree eq 0" \
  "add-labels-to-selected outdegsel" > outdeg.dot
\end{verbatim}
These directed-aware predicates complement the total-degree selector when working with asymmetric graphs.

\subsection{SVG Output}
\begin{verbatim}
node gruggle.js --ortho -i sample1.dot chain \
  "add-nodes 3 extra" \
  "to-svg sample1.svg"
\end{verbatim}
The SVG retains the light-blue, rounded-box styling and orthogonal edges when \texttt{--ortho} is used.

\end{document}
